% Labels of the draft are printed as \dref{label} (texttt), since this note is compiled on its own.
\providecommand{\dref}[1]{\texttt{\detokenize{#1}}}
\section{Referee 3, Major 2: the dilute--dilute limit against $k_T$ factorisation}\label{sec:ref3dilute}
% Scripts: notes/scripts/ref3_dilute/d1_qmrk_integrand.py (+ d1.out), d2_polarisation_and_convolution.py (+ d2.out).
% Literature sources (arXiv sources read in this session, export.arxiv.org): hep-ph/0608154 (BSV),
% hep-ph/9802290 (FL), hep-ph/9810215 (DDS).

\noindent\textbf{Result in one paragraph.} In the dilute--dilute limit (one gluon exchanged with the target in the amplitude and one in the conjugate amplitude, projectile two-point function $\langle\rho\rho\rangle$) the \emph{real} part of the NLO result, Groups~II and~III, i.e.\ $\int dp^+\!\int_v|\mathcal A^{ab}_{ij}|^2$ with the two-gluon amplitude of Eq.~(\dref{eq:A2gamp}) of the draft, is \textbf{exactly} the Fadin--Lipatov quasi-multi-Regge (QMRK) two-gluon production cross section, in the form $A_{\rm gluons}$ printed by Bartels, Sabio Vera and Schwennsen (BSV, hep-ph/0608154, Eq.~\texttt{eq:Agluons}), with the standard BFKL normalisation $\bar\alpha_s/\pi$ and no free constant: the ratio is $1$ to $2\times10^{-15}$ at random kinematics, for all momentum fractions $0<x<1$, for $SU(2)$ and $SU(3)$, and also for a transverse polarisation dimension $d_\perp\neq2$ against BSV's $(D-2)$ terms (CDR). Dropping any piece of the amplitude destroys the agreement (by factors $0.2$--$10$). The agreement requires the commutator (contact) parts of the two-gluon amplitude generated by the independent-emission terms ($\mf B_2$ with $\{\rho,\rho\}$, $\mathcal A^{(1,2)}_{AA}$); the anticommutator parts cancel exactly at first order in the target field. Both strong-ordering limits of the verified integrand give the real part of the LO BFKL kernel, on the projectile side ($p^+\gg k^+$) and on the target side ($p^+\ll k^+$), with $\bar\alpha_s=\alpha_sN_c/\pi$; the collinear limit is that of BSV. The finite real remainder in the paper's plus-momentum scheme is \emph{identical} to BSV's QMRK-minus-MRK remainder ($dp^+/p^+=-dy_2$ at fixed $k^+$; checked numerically): for the real emission no scheme conversion is needed. The \emph{virtual} part (Group~I) was reduced to the amplitude level in the dilute limit (Sec.~\ref{sec:d-virt}); its loop integrals were \textbf{not} done and it was \textbf{not} compared with the one-loop Lipatov vertex: this part of the referee's test remains open. Row~I.II alone has the rapidity structure $-\omega_B(\v K)\ln\frac{k^+}{\Lambda}$ with $\v K=\v k,\v q_a$ instead of the Regge $\omega(\v q_b)$; no amplitude-level cancellation of the $f^{abc}\{\rho^b,\rho^c\}$ part of $\mf A^{(3)}$ was found (it vanishes identically for a point-like projectile). For the constant the verified real part gives a sharp \emph{condition}, Eq.~(\ref{d:cond}): since $K^{(v)}_r$ has no $\frac{67}{9}$ and the real part is FL's, Group~I must contain no net $\frac{\alpha_sN_c}{2\pi}(\frac{67}{9}-\frac{\pi^2}{3})$ beyond scheme terms; with Row~I.II as the only LO-type constant (the draft's statement) the paper would exceed FL in the momentum-weighted inclusive constant by $4.155\,N_c\frac{\alpha_s}{2\pi}$ (or $(\frac{37}{4}-\pi^2)N_c\frac{\alpha_s}{2\pi}$ with the other fragmentation convention). This is an unresolved difference, conditional on the unevaluated Rows~I.III, I.V and contact terms; it contradicts the draft's explanation that the other half of the constant sits in the (real) collinear terms. The coefficient of $\ln\mu^2$ and the scale of the coupling agree. No disagreement was established. This check is made at the level of the amplitudes and coefficients of the paper (Eqs.~(\dref{eq:A2gamp}), (\dref{eq:B2mom}), (\dref{eq:C1}), (\dref{eq:A1})); it does \emph{not} test the coordinate-space finite terms as printed (Appendix~\dref{app:finite}), which would require their dilute limit to be evaluated and compared with the momentum-space formula of Sec.~\ref{sec:d-assembled}.

\medskip
\noindent\textbf{Sources read in this session} (arXiv sources from export.arxiv.org, kept in the scratchpad): BSV hep-ph/0608154 (Eqs.~\texttt{eq:Agluons}, \texttt{eq:defsigma}, \texttt{eq:e}, \texttt{eq:kernelv}, \texttt{eq:kernelagluons}, \texttt{eq:defbsquare}, \texttt{eq:jetvertexnloy}, \texttt{eq:sajs0lo}, \texttt{eq:runningcoupling}, Sec.~3 and the identity $2\v q_a\cdot\v k_1-\v q_a^2=\hat t+\v k_1^2/x$ of their Appendix, which fixes $\hat t$); Fadin--Lipatov hep-ph/9802290 (FL; Eqs.~(5)--(7)); Del Duca--Schmidt hep-ph/9810215 (DDS; Eq.~\texttt{looplip}). Not read in this session and therefore not used: the one-loop vertex papers of Fadin, Fiore et al.\ and Fadin--Fucilla--Papa, the Fadin--Lipatov QMRK paper hep-ph/9602287, and the Chirilli--Xiao--Yuan papers (their overlap with this calculation, the final-state collinear factorisation, was compared in \texttt{lit\_dglap\_rc.tex}; a dilute-limit comparison with the hybrid formulas was not done, since there the projectile is a collinear parton and the observable is forward).

\subsection{The limit and the counting}\label{sec:d-limit}
Write $U^{ab}(\v x)=\delta^{ab}+u^{ab}(\v x)$, $u^{ab}=f^{acb}\alpha^c$ (the convention below Eq.~(\dref{eq:ALOmom}); $U=e^{ig\alpha^aT^a}$ up to $\alpha\to-\alpha$ and the factor $g$ absorbed), $U^\dagger=U^T=1-u+O(u^2)$, and $\langle\alpha^c(\v l)\alpha^d(-\v l')\rangle=\delta^{cd}(2\pi)^2\delta^{(2)}(\v l-\v l')T(\v l)$, $\langle\rho^b(\v q')\rho^{b'}(-\v q)\rangle=\delta^{bb'}(2\pi)^2\delta^{(2)}(\v q-\v q')P(\v q)$ as in Eq.~(\dref{eq:kTfact}).
\begin{itemize}
\item \emph{Target.} Every amplitude vanishes at $U=1$, so each of $\mathcal A$, $\mathcal A^\dagger$ starts at $O(u)$. The dilute limit keeps $O(u)$ in the amplitude and in the conjugate amplitude: two-gluon exchange with the target in amplitude$\times$c.c., $\propto\langle\alpha\alpha\rangle$. Terms $O(u^2)\times O(u)$ are three-gluon (C-odd) exchanges and are dropped.
\item \emph{Projectile.} One $t$-channel gluon in the amplitude and one in the c.c.: terms with $\langle\rho\rho\rangle$. Since the charges do not commute, products of two or three $\rho$'s contribute to this order through $[\rho^a(\v x),\rho^b(\v y)]=if^{abc}\rho^c(\v x)\delta^{(2)}(\v x-\v y)$ (contact terms). The symmetrised three-point function $o$ is dropped (odderon exchange with the projectile; by the draft, Section~\dref{sec:checks}, its contribution starts with three gluons exchanged with the target, a statement not re-checked here).
\item \emph{Order.} LO is $O(g^2\langle\rho\rho\rangle\langle\alpha\alpha\rangle)$, the NLO terms considered are $O(g^4\langle\rho\rho\rangle\langle\alpha\alpha\rangle)$, i.e.\ $O(\alpha_s)$ relative to LO. For translation-invariant correlators everything is diagonal in momentum space.
\end{itemize}
Notation: $\int_{\v q}\equiv\int\frac{d^2q}{(2\pi)^2}$, $\bar\alpha_s=\frac{\alpha_sN_c}{\pi}$,
\begin{align}
\mathcal N\equiv\frac{\alpha_sN_c(N_c^2-1)}{\pi^2}\,\frac{S_\perp}{k^+}\,,\qquad
\phi_P(\v q)\equiv\frac{P(\v q)}{\v q^2}\,,\qquad \phi_T(\v l)\equiv\v l^2\,T(\v l)\,.\label{d:not}
\end{align}
$\phi_P$ and $\phi_T$ are the unintegrated gluon distributions (both $\propto\rho^2/\v q^2$; the target correlator $T$ is that of the field $\alpha\sim\rho_T/\v l^2$).

\subsection{Leading order}
Eq.~(\dref{eq:kTfact}) of the draft is, with $\frac{(\v k-\v q)^2}{\v k^2\v q^2}P(\v q)T(\v k-\v q)=\frac{\phi_P(\v q)\phi_T(\v k-\v q)}{\v k^2}$,
\begin{align}
\frac{d^3N_{\rm LO}}{d^3k}=\mathcal N\int_{\v q_a}\phi_P(\v q_a)\,\phi_T(\v q_b)\,\frac{A^{(0)}(\v q_a,\v q_b)}{\v q_a^2\,\v q_b^2}\,,\qquad A^{(0)}\equiv\frac{\v q_a^2\v q_b^2}{\boldsymbol\Delta^2}\,,\quad\v q_b=\v k-\v q_a\,,\quad\boldsymbol\Delta=\v q_a+\v q_b=\v k\,,\label{d:LO}
\end{align}
the $k_T$-factorised form with the squared Lipatov vertex $A^{(0)}$, in BSV variables ($\v q_a$ from the projectile, $\v q_b$ from the target, $\boldsymbol\Delta$ the produced momentum). In the LO amplitude, $\mathcal A^a_{{\rm LO}i}=\frac{\sqrt2g}{\sqrt{k^+}}\int_{\v l}u^{ab}(\v l)\rho^b(-\v q_a)V^{(0)}_i$, $V^{(0)}=\frac{\v k}{\v k^2}-\frac{\v q_a}{\v q_a^2}$, $\v q_a=\v k+\v l$ (Eq.~(\dref{eq:ALOmom}); re-derived here from Eq.~(\dref{eq:A1}) and the Fourier conventions).

\subsection{Real emission: Groups II and III}\label{sec:d-real}
\paragraph{The two-gluon amplitude at $O(u)$.} Insert $U=1+u$ in Eq.~(\dref{eq:A2gamp}); $\bar{\mf A}=\mf A+\delta\mf A$, $\delta\mf A\equiv\mf A[\rho\to u\rho]$, and $L^a_i(k^+,\v w)\equiv u^{ac}(\v w)\mf A^c_i(k^+,-\v w)-\delta\mf A^a_i(k^+,-\v w)$ (its Fourier transform is the LO amplitude). At $O(u^0)$ every term vanishes. At $O(u)$:
\begin{align}
\mathcal A_B&\to u^{ac}(\v w)\,\mf B^{bc}_{2ji}+u^{bc}(\v v)\,\mf B^{ca}_{2ji}-\mf B^{ba}_{2ji}[\rho\to u\rho]\,,\notag\\
\mathcal A^{(1)}_{AA}&\to-\tfrac12\,\mf A^b_j(p^+,-\v v)\,L^a_i(k^+,\v w)\,,\qquad
\mathcal A^{(2)}_{AA}\to-\tfrac12\,\mf A^a_i(k^+,-\v w)\,L^b_j(p^+,\v v)\,,\notag\\
\mathcal A_{CA}&\to\int_{\v z}\mf C^{bac}_{1jik}(p^++k^+,\v z-\v w;p^+,\v w-\v v)\,L^c_k(p^++k^+,\v z)\,,\label{d:Aexp}
\end{align}
(the terms with $u(\v v)$ in $\mathcal A^{(1)}_{AA}$ and with $\delta\mf A(\v v)$ cancel).
 The anticommutator part of $\mf B_2$, Eq.~(\dref{eq:B2}), is $\mf B^{(2)bc}_{2ji}=\frac14\{\mf A^c_i(k^+,-\v w),\mf A^b_j(p^+,-\v v)\}$, so that
\begin{align}
\mathcal A_B^{(2)}+\mathcal A^{(1)}_{AA}+\mathcal A^{(2)}_{AA}=\tfrac14\big[L^a_i(k^+,\v w),\mf A^b_j(p^+,-\v v)\big]+\tfrac14\big[L^b_j(p^+,\v v),\mf A^a_i(k^+,-\v w)\big]+O(u^2)\,.\label{d:comm}
\end{align}
\emph{The anticommutators cancel exactly at $O(u)$}; what is left is linear in $\rho$ by $[\rho,\rho]=if\rho\delta$ (for commuting charges it vanishes, as it must for a coherent state). Hence the two-gluon amplitude at $O(u)$ is, as an operator, linear in $\rho$: real emission in the dilute limit is a functional of $\langle\rho\rho\rangle$ only.

\paragraph{Momentum space.} With observed gluon $(k^+,\v k)$ (colour $a$, polarisation $i$), second gluon $(p^+,\v p)$ ($b$, $j$), $P^+=k^++p^+$, $x=k^+/P^+$, target momentum $-\v l$, projectile momentum $\v q_a=\v k+\v p+\v l$, the Fourier transform $\int_{w,v}e^{-i\v k\cdot\v w-i\v p\cdot\v v}$ of the bracket of Eq.~(\dref{eq:A2gamp}) is $\int_{\v l}\alpha^e(\v l)\rho^d(-\v q_a)\,M^{abed}_{ji}$ with
\begin{align}
M^{abed}_{ji}=\ &f^{aec}f^{cbd}\,\beta_{ji}(\v k+\v l,\v p)+f^{bec}f^{acd}\,\beta_{ji}(\v k,\v p+\v l)-f^{abc}f^{ced}\,\beta_{ji}(\v k,\v p)\notag\\
&+\frac{ig^2}{2\sqrt{k^+p^+}}\Big\{f^{aec}f^{cbd}\Big[\frac{(k+l)^ip^j}{(\v k+\v l)^2\v p^2}-\frac{k^ip^j}{\v k^2\v p^2}\Big]+f^{bec}f^{cad}\Big[\frac{k^i(p+l)^j}{\v k^2(\v p+\v l)^2}-\frac{k^ip^j}{\v k^2\v p^2}\Big]\Big\}\notag\\
&-\frac{ig^2\sqrt{k^+p^+}}{(P^+)^2}\,f^{cba}f^{ced}\,\frac{1}{\boldsymbol\kappa^2}\Big[(\boldsymbol\kappa\cdot\v V)\delta_{ji}-\frac{P^+}{p^+}\kappa^jV^i-\frac{P^+}{k^+}\kappa^iV^j\Big],\label{d:M}
\end{align}
$\boldsymbol\kappa=(p^+\v k-k^+\v p)/P^+$, $\v V=\frac{\v k+\v p}{(\v k+\v p)^2}-\frac{\v q_a}{\v q_a^2}$ (the LO vertex of the parent), and $f^{abc}\rho^c(-\v k-\v p)\beta_{ji}(\v k,\v p)$ the one-$\rho$ part of $\mf B^{ba}_{2ji}(k^+,\v k;p^+,\v p)$, Eq.~(\dref{eq:B2mom}) (three terms: commutator part of the double emission, triple-gluon vertex, instantaneous emission). The first line is $\mathcal A_B^{(1)}$, the second Eq.~(\ref{d:comm}), the third $-\mathcal A_{CA}$ (the $\delta$-function of $\mf C_1$, Eq.~(\dref{eq:C1}), makes the pair integral $\int_{\v r}e^{-i\boldsymbol\kappa\cdot\v r}\Kv(\v r)=-2\pi i\boldsymbol\kappa/\boldsymbol\kappa^2$).

\paragraph{Identity (verified numerically; Eq.~(\ref{d:comm}) and Eq.~(\ref{d:M}) are derived by hand from the draft's formulas and enter only through it).}
\begin{align}
\sum_{a,b,e,d}\sum_{i,j}\big|M^{abed}_{ji}\big|^2=\frac{g^4\,N_c^2(N_c^2-1)}{k^+p^+\,\v q_a^4}\;A_{\rm gluons}(\v q_a,\v q_b;\v k_1=\v k,\v k_2=\v p;x)\Big|_{D=d_\perp+2}\,,\qquad\v q_b=-\v l\,,\label{d:identity}
\end{align}
with $A_{\rm gluons}$ of BSV Eq.~\texttt{eq:Agluons}, $\boldsymbol\Lambda=(1-x)\v k_1-x\v k_2$, $\hat s=\boldsymbol\Lambda^2/(x(1-x))$, $\Sigma=\hat s+\boldsymbol\Delta^2$, $E$ of BSV Eq.~\texttt{eq:e}, and $\hat t=2\v q_a\cdot\v k_1-\v q_a^2-\v k_1^2/x$, $\hat u=2\v q_a\cdot\v k_2-\v q_a^2-\v k_2^2/(1-x)$ (BSV Appendix; $\hat t=(q_a-k_1)^2$ in Minkowski signature). Script \texttt{d1\_qmrk\_integrand.py}, output \texttt{d1.out}, \texttt{d2.out}:
\begin{center}
\renewcommand{\arraystretch}{1.2}
\begin{tabular}{|p{0.52\textwidth}|p{0.4\textwidth}|}\hline
check & result\\\hline
ratio of the two sides of Eq.~(\ref{d:identity}), 12 random points each, $0.05<x<0.95$, $SU(2)$ / $SU(3)$ & $1$, max.\ deviation $1.3\times10^{-15}$ / $1.9\times10^{-15}$\\
$d_\perp=1.7,\,3,\,6$ (term $\propto\delta_{ji}$ traced in $d_\perp$ dimensions) vs.\ BSV at $D=d_\perp+2$ & $1$ to $10^{-14}$ (5 points)\\
$x\to0$ and $x\to1$ at fixed transverse momenta & both sides $\to\v q_a^2\v q_b^2/(\v k^2\v p^2)$ (at $x=10^{-5}$, $1-10^{-5}$: ratio $1\pm3\times10^{-7}$)\\
without the commutator terms (2nd line of Eq.~(\ref{d:M})) & ratio $0.18$--$4.1$\\
commutator terms doubled / without $\mathcal A_{CA}$ / without $\mf B_2[u\rho]$ & $1.0$--$5.2$ / $0.25$--$5.1$ / $0.90$--$9.7$\\
convolution with $\phi_{P,T}=\v q^2e^{-\v q^2/Q^2_{P,T}}$, $0.1<x<0.9$, $|\v p|>0.25$, $|\v k|=0.7,1.5,3$ (same MC points) & equal to $2.4\times10^{-15}$\\\hline
\end{tabular}
\end{center}
The normalisation constant is exactly one; nothing was fitted. Since in the strong-ordering limits $A_{\rm gluons}\to\v q_a^2\v q_b^2/(\v k_1^2\v k_2^2)$ (product of two Lipatov vertices; BSV: $A_{\rm MRK}=-4A_{(4)}$), Eq.~(\ref{d:identity}) also fixes the overall normalisation of the real emission against the paper's own LO.

\paragraph{Result.} With $dN/d^3k=(2\pi)^{-3}\int dp^+\!\int\frac{d^2p}{(2\pi)^2}|\tilde{\mathcal A}|^2$, $\tilde{\mathcal A}=\frac{2}{\sqrt{2\pi}}\int_{\v l}\alpha\rho\,M$, Eq.~(\dref{eq:N_amplitudes}),
\begin{align}
\boxed{\;\frac{d^3N^{\rm(II+III)}}{d^3k}\Big|_{\rm dilute}=\mathcal N\int_{\v q_a}\int_{\Lambda}^{\vee}\frac{dp^+}{p^+}\int d^2p\;\frac{\bar\alpha_s}{\pi}\;\phi_P(\v q_a)\,\phi_T(\v q_b)\,\frac{A_{\rm gluons}(\v q_a,\v q_b;\v k,\v p;x)}{\v q_a^2\,\v q_b^2}\,,\qquad\v q_b=\v k+\v p-\v q_a\,,\quad x=\frac{k^+}{k^++p^+}\,.\;}\label{d:real}
\end{align}
Longitudinal ranges: the one-$\rho$ part of $\mf B_2$ and the parent of $\mf C_1$ require $k^++p^+<\vee$, the commutator terms only $p^+<\vee$; the difference is $O(k^+/\vee)$. In $d_\perp=2-\epsilon$: $d^2p\to\mu^\epsilon d^{2-\epsilon}p$ and $A_{\rm gluons}$ at $D=4-\epsilon$ (BSV write $D=4+2\epsilon$). Each produced gluon is counted once (single-inclusive number density); BSV's factor $2$ in configuration (c) of their jet vertex is the same statement for a jet.

\paragraph{Consequences.}
\begin{enumerate}
\item \emph{Rapidity logarithms, real parts.} Both limits $x\to0$ ($p^+\gg k^+$) and $x\to1$ ($p^+\ll k^+$) give the same function $\v q_a^2\v q_b^2/(\v k^2\v p^2)$. With $\v q'=\v q_a-\v p$ (projectile side) it is $\frac{\bar\alpha_s}{\pi}\frac{d^2p}{\v p^2}\phi_P(\v q'+\v p)\times\phi_T(\v k-\v q')/\v k^2$, with $\v q'_b=\v k-\v q_a$ (target side) $\phi_P(\v q_a)/\v k^2\times\frac{\bar\alpha_s}{\pi}\frac{d^2p}{\v p^2}\phi_T(\v q'_b+\v p)$: the real parts of one step of LO BFKL evolution of $\phi_P$ over $\ln\frac{\vee}{k^+}$ and of $\phi_T$ over $\ln\frac{k^+}{\Lambda}$, with $\bar\alpha_s=\alpha_sN_c/\pi$. Hence the finite real terms in the paper's scheme are
\begin{align}
\frac{d^3N^{\rm(II+III)}}{d^3k}\Big|_{\rm dilute}=\mathcal N\int_{\v q_a}\!\int d^2p\,\frac{\bar\alpha_s}{\pi}\,\phi_P\phi_T\Big\{\ln\frac{\vee}{\Lambda}\,\frac{1}{\v k^2\v p^2}+\int_0^1\frac{dx}{x(1-x)}\Big[\frac{A_{\rm gluons}}{\v q_a^2\v q_b^2}-\frac{1}{\v k^2\v p^2}\Big]\Big\}+O\Big(\frac{k^+}{\vee},\frac{\Lambda}{k^+}\Big),\label{d:realsplit}
\end{align}
($\int_\Lambda^\vee dp^+/p^+=\int dx/(x(1-x))$ over $\frac{k^+}{k^++\vee}<x<\frac{k^+}{k^++\Lambda}$); the curly bracket is IR divergent (soft $\v p\to0$, collinear $\boldsymbol\Lambda\to0$) and becomes finite only with Group~I and the fragmentation counterterm, in dimensional regularisation.
\item \emph{Collinear region and polarisation scheme.} The collinear terms $A_{(1)}+A_{(2)}$ of BSV (their Eq.~\texttt{eq:agluonscollinear}) are those of the paper's real emission. With the polarisation trace $\delta_{ii}=d_\perp$ the paper's two-gluon amplitude reproduces BSV's $(D-2)$ dependence for any $d_\perp$, i.e.\ the amplitude has the CDR structure. Whether the finite terms of the paper are evaluated with $\delta_{ii}=d_\perp$ throughout (the referee's Major~4) is \emph{not} decided by this check; it only shows that with this choice the real emission is the CDR one of FL/BSV.
\item \emph{Contact terms (Major 1).} The commutator terms of Eq.~(\ref{d:comm}) come from $\mf B_2^{(2)}$ and $\mathcal A_{AA}^{(1,2)}$, i.e.\ from the rows with three and four charges (II.I, II.II, II.III, III.I, III.II, III.IV, III.V, III.VI). In the cross section they give $2{\rm Re}[(\mathcal A_B^{(1)}-\mathcal A_{CA})^\dagger C]+|C|^2$, $C$ the commutator terms. $|C|^2\propto1/p^+$ is purely logarithmic (no finite part, consistent with the statement that the four-charge rows contribute only to the evolution terms); the interference is not. Without these two-point contact contributions the QMRK result is \emph{not} reproduced (table). So in the real sector the contact terms of the cubic rows cannot cancel among themselves: they are required. Whether the contact terms printed in the colour-singlet section of the paper equal the dilute-limit expressions above was not checked.
\end{enumerate}

\subsection{Rapidity logarithms and the scheme}\label{sec:d-logs}
In the dilute target limit the BFKL and JIMWLK terms are one step of linear evolution of the projectile and of the target unintegrated gluon distributions:
\begin{align}
\frac{d^3N_{\rm BFKL}}{d^3k}+\frac{d^3N_{\rm JIMWLK}}{d^3k}\Big|_{\rm dilute}&=\mathcal N\int_{\v q}\frac{1}{\v k^2}\Big[\ln\frac{\vee}{k^+}\,(\mathcal K\otimes\phi_P)(\v q)\,\phi_T(\v k-\v q)+\ln\frac{k^+}{\Lambda}\,\phi_P(\v q)\,(\mathcal K\otimes\phi_T)(\v k-\v q)\Big],\notag\\
(\mathcal K\otimes\phi)(\v q)&=\frac{\bar\alpha_s}{\pi}\int\frac{d^2p}{\v p^2}\Big[\phi(\v q+\v p)-\frac{\v q^2}{\v p^2+(\v q+\v p)^2}\,\phi(\v q)\Big].\label{d:BFKL}
\end{align}
\emph{What is established:} the real parts and their normalisation follow from the verified integrand, item~1 above; the virtual (trajectory) parts rest on the paper's own checks (Section~\dref{sec:check_BFKL}: exact one-gluon evolution of a dipole at finite $N_c$; Section~\dref{sec:check_JIMWLK}: $-H_{\rm JIMWLK}$ on $\mathcal O_{\rm LO}$), whose linearisation is the standard LO BFKL equation; the momentum-space trajectory was not re-derived here.

\emph{Scheme.} Both logarithms refer to cut-offs on the plus momentum of the soft gluon ($\Lambda<p^+<\vee$). BFKL/$k_T$ factorisation uses rapidity or energy variables with a scale $s_0$. Since $y_p-y_k=\ln\frac{p^+}{k^+}+\frac12\ln\frac{\v k^2}{\v p^2}$, the paper's window is, at transverse momentum $\v p$, the rapidity window of width $\ln\frac{\vee}{\Lambda}$ shifted by $-\frac12\ln\frac{\v p^2}{\v k^2}$. On the projectile side the paper's variable $\ln\frac{\vee}{k^+}$ is the plus-momentum fraction, as in BSV's asymmetric choice $s_0=\v k_J^2$, $s_{AJ}/s_0=1/\alpha$ (their Eq.~\texttt{eq:sajs0lo}); on the target side BSV use the minus-momentum fraction $1/\beta$, while the paper uses $\ln\frac{k^+}{\Lambda}$; for an emission at $\v p$, $\ln\frac{k^+}{p^+}=\ln\frac{p^-}{k^-}+\ln\frac{\v k^2}{\v p^2}$. The two schemes therefore differ by terms $\frac{\bar\alpha_s}{\pi}\int d^2p\,\ln\frac{\v p^2}{\v k^2}\times$(target-side real kernel minus the corresponding virtual term) acting on $\phi_T$ (and in a symmetric $s_0=|\v q_a||\v q_b|$ choice by the analogous terms on both sides); these are of the ``energy-scale'' type, not pure numbers times LO. In the sum of the real terms, Eq.~(\ref{d:realsplit}), the shift cancels because both strong-ordering limits are the same function; it does not cancel in the split between projectile and target evolution and in the virtual terms. \textbf{Exact result for the real part (script \texttt{d3\_real\_scheme.py}, \texttt{d3.out}).} At fixed $k^+$, $\frac{dp^+}{p^+}=\frac{dx}{x(1-x)}=-dy_2$ \emph{exactly} ($y_2$ the rapidity of the second gluon relative to the observed one, $x=|\v k|/(|\v k|+|\v p|e^{y_2})$, BSV's definition). Hence the paper's finite real remainder, the $x$ integral of Eq.~(\ref{d:realsplit}), is identical to BSV's QMRK-minus-MRK remainder, Eq.~\texttt{eq:defbsquare}, in the limit $s_\Lambda\to\infty$:
\begin{align}
\int_0^1\frac{dx}{x(1-x)}\Big[\frac{A_{\rm gluons}}{\v q_a^2\v q_b^2}-\frac{1}{\v k^2\v p^2}\Big]=\lim_{s_\Lambda\to\infty}\Big\{\int dy_2\,\frac{A_{\rm gluons}\,\theta(s_\Lambda-s_{12})}{\v q_a^2\v q_b^2}-\frac{1}{\v k^2\v p^2}\ln\frac{s_\Lambda^2}{\v k^2\v p^2}\Big\}\label{d:schemeid}
\end{align}
(single-inclusive normalisation; BSV's factor $\frac12$ in front of their MRK term belongs to their pair normalisation, cf.\ their factor $2$ in configuration (c)). Numerically, at three random points: $-2.62950455$, $-0.66773731$, $-6.13068770$ (paper) against $-2.62950704$, $-0.66773773$, $-6.13068408$ (BSV form at $s_\Lambda=10^6$; at $10^2$, $10^4$ the differences fall like $1/s_\Lambda$). \emph{So for the real emission there is no conversion term at all}: the plus-momentum window $\Lambda<p^+<\vee$ has, at every $\v p$, the rapidity width $\ln\frac{\vee}{\Lambda}$, it is only shifted by $-\frac12\ln\frac{\v p^2}{\v k^2}$, and since both strong-ordering limits are the same function the shift cancels between the two ends. The scheme dependence is entirely in (i) the identification of $\ln\frac{\vee}{k^+}$, $\ln\frac{k^+}{\Lambda}$ with BSV's $\ln\frac{s_{AJ}}{s_0}$, $\ln\frac{s_{BJ}}{s_0'}$ (i.e.\ in the definition of the evolved $\phi_{P,T}$), and (ii) the virtual trajectory terms, where the split of the loop-gluon window at $p^+=k^+$ (paper) instead of $y_p=y_k$ (Regge) moves $\frac12\ln\frac{\v p^2}{\v k^2}\times[\text{integrand of }\omega(\v q_a)-\text{integrand of }\omega(\v q_b)]$ between the two sides. Item (ii) requires the paper's dilute Group~I (Sec.~\ref{sec:d-virt}) and was not evaluated.

\subsection{Coupling, collinear terms, constant}\label{sec:d-const}
\paragraph{$\beta_0$ and $\mu$ (agrees).} In the local form of Final\_terms\_fixed (Eqs.~\texttt{fx-locsplit}, \texttt{fx-equiv}), whose dilute limit is immediate in momentum space ($\Pi(\xi)=\xi^{-2}\widehat L(\v k/\xi)$ is the LO spectrum of the parent):
\begin{align}
\frac{d^3N_{\rm LO}}{d^3k}+\frac{d^3N^{\rm loc}_{\rm RC}}{d^3k}&=\frac{d^3N_{\rm LO}}{d^3k}\Big[1-\frac{\alpha_s\beta_0}{4\pi}\ln\frac{\v k^2}{\mu^2}\Big],\qquad
\frac{d^3N^{\rm loc}_{\rm DGLAP}}{d^3k}=\frac{\alpha_s}{2\pi}\,\mathcal N_0\ln\frac{\v k^2}{\mu^2}\int_0^1d\xi\,\mathcal P_{gg}(\xi)\,\Pi(\xi)\,,\label{d:RC}
\end{align}
plus the non-local, $\mu$-independent remainder $N_{\mathcal R}$. FL Eq.~(7): $K^{(1)}_r\supset K^B_r\,\bar g^2_\mu\frac{11}{3}\frac1\epsilon\big(1-(\boldsymbol\Delta^2/\mu^2)^\epsilon\big)\to-\frac{\alpha_sN_c}{4\pi}\frac{11}{3}\ln\frac{\boldsymbol\Delta^2}{\mu^2}K^B_r$, $\boldsymbol\Delta=\v k$; BSV Eq.~\texttt{eq:runningcoupling}: $\mathcal V^{(B)}(1-\frac{\bar\alpha_s\beta_0}{4N_c}\ln\frac{\v k_J^2}{\mu^2})$ (printed with $\alpha_s$; consistent with their finite terms only for $\bar\alpha_s$, see \texttt{lit\_finite\_dilute.tex}). Same coefficient $\frac{\alpha_s\beta_0}{4\pi}$, same scale (the produced momentum). The $\mathcal P_{gg}\ln\mu^2$ term is the $\overline{\rm MS}$ fragmentation counterpart, absent for a jet.
\paragraph{Constant $\frac{67}{9}-\frac{\pi^2}{3}$: a consistency condition (not decided).} In the paper the LO-type constant sits in the virtual part: Row~I.II, one-$\rho$, $c(\v K)$ of Sec.~\ref{sec:d-virt}, whose constant gives $+\frac{\alpha_sN_c}{2\pi}(\frac{67}{9}-\frac{\pi^2}{3})\times$LO $\simeq1.98\,\alpha_s\times$LO ($N_c=3$). In BFKL the one-loop vertex contains no $\frac{67}{9}$: DDS Eq.~\texttt{looplip} ($c_\Gamma\{N_c[\dots-\frac{1}{\epsilon^2}(\frac{\mu^2}{\v k^2})^\epsilon-\frac12\ln^2\frac{t_1}{t_2}+\frac{\pi^2}{3}]-\frac{\beta_0}{2\epsilon}-\dots\}$, depending on the Reggeisation scale $\tau$) and BSV Eq.~\texttt{eq:kernelv} ($K^{(v)}_r$: $-\frac{2}{\epsilon^2}$, $\pi^2$, $\frac{\beta_0}{N_c\epsilon}$, $\ln^2\frac{\v q_a^2}{\v q_b^2}$, no $\frac{67}{9}$); in FL Eq.~(7) the term $\frac{\alpha_sN_c}{4\pi}(\frac{67}{9}-\frac{\pi^2}{3})$ appears only after the real QMRK (and $q\bar q$) production is integrated over the pair at fixed $\boldsymbol\Delta$.

\emph{Consequence of the verified real part.} The paper's real emission is the QMRK cross section, Eq.~(\ref{d:identity}), with the same polarisation structure (CDR) and transverse dimensional regularisation, and its finite remainder coincides identically with BSV's, Eq.~(\ref{d:schemeid}). Therefore, for any pair-inclusive quantity built from the paper's amplitudes (virtual at $\boldsymbol\Delta$ plus the two-gluon production integrated at fixed $\boldsymbol\Delta$), the real part contributes exactly FL's $\frac{\alpha_sN_c}{4\pi}(\frac{67}{9}-\frac{\pi^2}{3})$, provided the real rows are integrated with $\delta_{ii}=d_\perp$ (Major~4). The rapidity-scheme terms (item (ii) of Sec.~\ref{sec:d-logs}) are $\ln(\v p^2/\v k^2)$-weighted LO kernels and trajectories; they contain no moments of the splitting function and hence no $\frac{67}{9}$. Agreement with FL thus requires, up to these scheme terms, that the paper's dilute Group~I contains \emph{no net} $\frac{67}{9}-\frac{\pi^2}{3}$:
\begin{align}
\big[\text{Group I}\big]_{\rm dilute}\Big|_{c}=\underbrace{+\frac{\alpha_sN_c}{2\pi}\Big(\frac{67}{9}-\frac{\pi^2}{3}\Big)}_{\text{Row I.II}}+\big[\text{Rows I.III, I.V, contact terms}\big]_{c}\overset{!}{=}\text{(scheme terms)}\,.\label{d:cond}
\end{align}
Either Rows~I.III, I.V or the contact terms contain a compensating $-\frac{\alpha_sN_c}{2\pi}(\frac{67}{9}-\frac{\pi^2}{3})$-type term (the $\xi$-dependent energy denominators $D(\xi)=\xi(y-x)^2+(1-\xi)(y-z)^2$ of Row~I.V, Eq.~\texttt{fx-V2}, are the only place in Group~I where a $\int d\xi\,P(\xi)\ln(\xi(1-\xi))$-type moment could arise), or the paper's result differs from FL's by $+\frac{\alpha_sN_c}{2\pi}(\frac{67}{9}-\frac{\pi^2}{3})\times$LO in this quantity (up to $\pi^2$ terms from the scheme conversion). This is in tension with the draft's explanation (Section~\dref{sec:lit_constant}) that ``the other half must reside in the collinear finite terms'' (Rows~II.I, II.IV, III.III): those are real rows, and their sum is now fixed to be the QMRK cross section, whose pair-inclusive integral gives $+\frac{\alpha_sN_c}{4\pi}(\frac{67}{9}-\frac{\pi^2}{3})$, not $-\frac{\alpha_sN_c}{4\pi}(\dots)$. For the single-inclusive gluon the same statement holds for the momentum-weighted sum over both gluons, up to the fragmentation-scheme term of the $\xi^{-d}$ Jacobian, whose momentum moment is $\frac{\alpha_s}{2\pi}\int_0^1d\xi\,\xi\,\mathcal P_{gg}\ln\xi^2=\frac{\alpha_s}{2\pi}N_c\big(\frac{65}{36}-\frac{2\pi^2}{3}\big)$ ($=-14.3225\,\frac{\alpha_s}{2\pi}$ at $N_c=3$, as in \texttt{lit\_dglap\_rc}). Hence, \emph{if} Row~I.II carried the only LO-type constant of Group~I (the draft's statement), the momentum-weighted inclusive constant of the paper would exceed FL's by $\frac{\alpha_s}{2\pi}N_c(\frac{67}{9}-\frac{\pi^2}{3})=4.155\,N_c\frac{\alpha_s}{2\pi}$ if the paper's fragmentation function is the one obeying the momentum sum rule, or by $\frac{\alpha_s}{2\pi}N_c[(\frac{67}{9}-\frac{\pi^2}{3})+(\frac{65}{36}-\frac{2\pi^2}{3})]=\frac{\alpha_s}{2\pi}N_c(\frac{37}{4}-\pi^2)=-0.620\,N_c\frac{\alpha_s}{2\pi}$ if the CXY definition (Eq.~\texttt{fx-Delta} of Final\_terms\_fixed, $\Delta$ added) is the momentum-conserving one; neither is zero. Both numbers are conditional; the rapidity-scheme terms (ii) can add further $\pi^2$-type constants. \textbf{Which alternative holds is decided only by the dilute Group~I (Sec.~\ref{sec:d-virt}), which was not evaluated.} I report this as an unresolved difference, not as an error of the paper; I did not adjust anything.

\subsection{Virtual part: Group I in the dilute limit (reduced at amplitude level; loop integrals not done)}\label{sec:d-virt}
\paragraph{Reduction at $O(u)$.} From Eq.~(\dref{eq:A3amp}) with $U=1+u$, $\bar{\mf X}=\mf X+\delta\mf X$ ($\delta\mf X\equiv$ the part linear in $u$ of $\mf X[\rho\to U\rho]$), $L$ as in Sec.~\ref{sec:d-real} (derived by hand):
\begin{align}
\mathcal A_N+[\mathcal A_{A^{(3)}}]_{\rho\rho\rho}&\to\big\{\mf N_2,\,L^a_i(k^+,\v w)\big\}\quad(\text{Eq.~(\dref{eq:Nterms})}),\notag\\
[\mathcal A_{A^{(3)}}]_{\rho}&\to\frac{\sqrt2 g}{\sqrt{k^+}}\int_{\v l}u^{ab}(\v l)\rho^b(-\v q_a)\Big[c(\v k)\frac{\v k}{\v k^2}-c(\v q_a)\frac{\v q_a}{\v q_a^2}\Big]_i\quad(\text{after }\textstyle\int_we^{-i\v k\cdot\v w}),\notag\\
[\mathcal A_{A^{(3)}}]_{\rho\rho}&\to u^{ac}(\v w)\,[\mf A^{(3)c}_{i}]_{\rho\rho}-\delta[\mf A^{(3)a}_{i}]_{\rho\rho}\,,\qquad(f^{abc}\{\rho^b,\rho^c\}\ \text{structure, Eq.~(\dref{eq:A3tworho})}),\notag\\
\mathcal A_{B_3}&\to\int\frac{dp^+}{2\pi}\int_{\v x}\mf B^{\dagger ba}_{3ji}(k^+,\v w;p^+,-\v x)\,L^b_j(p^+,\v x)\,,\notag\\
\mathcal A_{AB_2}&\to2\int\frac{dp^+}{2\pi}\int_{\v x}\Big[u^{ad}(\v w)\mf A^{\dagger b}_j\mf B^{bd}_{2ji}+\mf A^{\dagger b}_ju^{bc}(\v x)\mf B^{ca}_{2ji}-\mf A^{\dagger b}_j\,\delta\mf B^{ba}_{2ji}\Big],\notag\\
\mathcal A_{CB_2}&\to2\int\frac{dp^+}{2\pi}\int_{\v x,\v z}\mf C^{\dagger}_1\big[(u(\v x)+u(\v z))\,\mf B_2-\delta\mf B_2\big]\,,\label{d:virtred}
\end{align}
with $c(\v K)=-\frac{\alpha_sN_c}{4\pi}\big\{[-\frac2\epsilon+\ln\frac{\v K^2}{\mu^2}][\frac{11}{3}-2\ln\frac{k^+}{\Lambda}+\ln(1-\frac{k^+}{\vee})]+\frac{\pi^2}{3}-\frac{67}{9}+\ln^2(1-\frac{k^+}{\vee})+\ln(1-\frac{k^+}{\vee})\ln\frac{\vee}{k^+}\big\}$ the ratio $[\mf A^{(3)}]_\rho/\mf A^{(1)}$ of Eq.~(\dref{eq:A3onerho}) (in $\mathcal A_{AB_2}$ the two terms $\pm\delta\mf A^\dagger\mf B_2$ cancel; arguments as in Eq.~(\dref{eq:A3amp})). Group~I is then $(2\pi)^{-3}\langle\mathcal A^\dagger_{\rm LO}\mathcal A_{(3)}\rangle+{\rm c.c.}$ with $\mathcal A_{\rm LO}\propto L$.
\begin{itemize}
\item \emph{Row I.I.} $\{\mf N_2,L\}$ is proportional to $\ln\frac{\vee}{\Lambda}$ with no other $p^+$ dependence: no finite part, consistent with the draft. For a point-like projectile ($\rho^a(\v x)=t^a\delta^{(2)}(\v x)$) $\mf N_2\propto\int d^{2-\epsilon}p/\v p^2=0$ in dimensional regularisation.
\item \emph{Row I.II, one-$\rho$.} Each term of the Lipatov vertex is multiplied by $1+c$: $c(\v k)$ for emission after, $c(\v q_a)$ for emission before the collision. Its rapidity-dependent part is $-\omega_B(\v K)\ln\frac{k^+}{\Lambda}$ exactly, with the LO trajectory of FL Eq.~(5), $\omega_B(\v K)=-\bar g^2_\mu(\frac{2}{\epsilon_{\rm FL}}+2\ln\frac{\v K^2}{\mu^2})=\frac{\alpha_sN_c}{2\pi}(\frac2\epsilon-\ln\frac{\v K^2}{\mu^2})$ ($d_\perp=2-\epsilon=2+2\epsilon_{\rm FL}$). Regge factorisation (BSV Eq.~\texttt{eq:news17nlo}) requires the target-side logarithm to multiply $\omega(\v q_b)$, $\v q_b=\v k-\v q_a$, for the whole vertex, and the projectile-side one to multiply $\omega(\v q_a)$; Row~I.II alone gives $\omega(\v k)$ and $\omega(\v q_a)$ on the two terms. The conversion must come from Rows~I.III, I.V and the contact terms; that it does is the content of the draft's JIMWLK and BFKL checks (Sections~\dref{sec:check_JIMWLK}, \dref{sec:check_BFKL}), not re-verified in momentum space.
\item \emph{The $f^{abc}\{\rho^b,\rho^c\}$ structures (question of an analogue of Eq.~(\ref{d:comm})).} In the real emission the cancellation is between products of single-emission coefficients ($\mf B_2^{(2)}=\frac14\{\mf A,\mf A\}$ and $\mathcal A_{AA}$), the coherent-state mechanism. In Group~I the anticommutator part of $\mf B_3$ is also of this type, $\mf B_3^{(2)}=\frac12\{\cdot,\cdot\}$ of two single-emission amplitudes (Eq.~(\dref{eq:B3}), third line), and so is $\mf B_2^{(2)}$ inside $\mathcal A_{AB_2}$, $\mathcal A_{CB_2}$; for these an analogous reduction to commutators is plausible but was not carried out. The $f^{abc}\{\rho^b,\rho^c\}$ part of $\mf A^{(3)}$ is \emph{not} a product of single emissions (two charges joined by the triple-gluon or instantaneous vertex) and does not vanish for commuting charges; at $O(u)$ in the amplitude it is therefore not removed by the coherent-state mechanism, and \textbf{no amplitude-level analogue of Eq.~(\ref{d:comm}) was found for it}. Its contribution at $O(u^2)\langle\rho\rho\rangle$ is through contact terms only (its symmetrised part is $\propto o$). It vanishes identically for a point-like projectile, $f^{abc}\{t^b,t^c\}=0$.
\item \emph{Point-like projectile.} For $\rho^a=t^a\delta^{(2)}(\v x)$ (a quark), Rows~I.I and the $\rho\rho$, $\rho\rho\rho$ parts of $\mf A^{(3)}$ vanish (above), and Group~I reduces to $[\mathcal A_{A^{(3)}}]_\rho$, $\mathcal A_{B_3}$, $\mathcal A_{AB_2}$ (finite part $=-$III.I) and $\mathcal A_{CB_2}$, with exact colour algebra ($t^at^b=\frac{1}{2N_c}\delta^{ab}+\frac12(d^{abc}+if^{abc})t^c$). This is the cleanest setting for the comparison with the one-loop vertex.
\end{itemize}
\paragraph{What was not done.} The loop integrals: $\int_\Lambda^\vee dp^+$ and $\int d^{2-\epsilon}p$ of $\mf B_3^\dagger L$ (principal-value denominators $p^+\v k^2-k^+\v p^2$, regions $p^+\gtrless k^+$), of $\mf A^\dagger\mf B_2$ and of $\mf C^\dagger_1\mf B_2$ (UV pole, cancelling the $\frac{11}{3}$ of Row~I.II according to the draft), followed by the subtraction of the BFKL/JIMWLK-list terms and the comparison with $K^{(v)}_r$ (BSV Eq.~\texttt{eq:kernelv}) / DDS Eq.~\texttt{looplip} after the trajectory conversion (ii) of Sec.~\ref{sec:d-logs}. The covariant one-loop vertex is available for arbitrary $D$ (Fadin, Fiore, Kotsky, Papa, hep-ph/0008006; Fadin, Fucilla, Papa, arXiv:2302.09868; titles seen in an arXiv listing, papers not read). This is a computation of the size of the original Group~I evaluation; it was not attempted in this session.

\subsection{The assembled momentum-space form}\label{sec:d-assembled}
Collecting Eqs.~(\ref{d:LO})--(\ref{d:RC}), in the dilute--dilute limit
\begin{align}
\frac{d^3N}{d^3k}\Big|_{\rm dilute}=\ &\mathcal N\int_{\v q_a}\phi_P(\v q_a)\phi_T(\v k-\v q_a)\,\frac{1}{\v k^2}\Big[1-\frac{\alpha_s\beta_0}{4\pi}\ln\frac{\v k^2}{\mu^2}\Big]+\frac{d^3N^{\rm loc}_{\rm DGLAP}}{d^3k}\notag\\
&+\ln\frac{\vee}{k^+}\,\mathcal N\!\int_{\v q}\frac{(\mathcal K\otimes\phi_P)(\v q)\phi_T(\v k-\v q)}{\v k^2}+\ln\frac{k^+}{\Lambda}\,\mathcal N\!\int_{\v q}\frac{\phi_P(\v q)(\mathcal K\otimes\phi_T)(\v k-\v q)}{\v k^2}\notag\\
&+\Big[\mathcal N\int_{\v q_a}\!\int d^2p\,\frac{\bar\alpha_s}{\pi}\phi_P\phi_T\int_0^1\frac{dx}{x(1-x)}\Big(\frac{A_{\rm gluons}}{\v q_a^2\v q_b^2}-\frac{1}{\v k^2\v p^2}\Big)+\frac{d^3N^{\rm(I)}_{\rm fin}}{d^3k}\Big|_{\rm dilute}\Big]_{\overline{\rm MS},\ \epsilon\to0}\,,\label{d:all}
\end{align}
where the last bracket is the sum of the finite terms of the paper, evaluated in $d_\perp=2-\epsilon$ with the $\overline{\rm MS}$ coupling and fragmentation counterterms taken at $\mu^2=\v k^2$ (so that the whole $\mu$ dependence is in the first line), then $\epsilon\to0$. Its real part is Eq.~(\ref{d:realsplit}) without the logarithm (verified integrand, Eq.~(\ref{d:identity})); it contains the collinear region of Row~II.IV, hence the non-local remainder $N_{\mathcal R}$ and the fragmentation-scheme terms. Its virtual part $N^{\rm(I)}_{\rm fin}$ (Group~I minus its BFKL/JIMWLK-list terms, including the LO-type constant $\frac{\alpha_sN_c}{2\pi}(\frac{67}{9}-\frac{\pi^2}{3})N_{\rm LO}$) was not reduced. $\mathcal N_0$ in Eq.~(\ref{d:RC}) is the LO prefactor of Final\_terms\_fixed; $\Pi(\xi)$ the LO spectrum of a parent of momentum $\v k/\xi$ in the same normalisation.

\subsection{Comparison table}
\begin{center}
\renewcommand{\arraystretch}{1.22}
\begin{tabular}{|p{0.30\textwidth}|p{0.16\textwidth}|p{0.46\textwidth}|}\hline
item & verdict & remark\\\hline
LO, dilute target & AGREES & $k_T$ factorisation with $|\text{Lipatov vertex}|^2$, Eq.~(\ref{d:LO})\\
real $gg$ production, integrand & AGREES exactly & paper (Groups II+III) $=$ BSV/FL QMRK $A_{\rm gluons}$, all $x$, $SU(2)$, $SU(3)$, $d_\perp$ arbitrary; ratio $1\pm2\times10^{-15}$\\
contact (commutator) terms of the real rows & required & without them QMRK fails (factors $0.2$--$4$)\\
rapidity logs, real parts & AGREES & $x\to0,1$: LO BFKL real kernel on $\phi_P$, $\phi_T$, $\bar\alpha_s/\pi$\\
rapidity logs, virtual parts & STRUCTURE & paper's own exact-LL checks; momentum-space trajectory not re-derived\\
rapidity scheme, real part & IDENTICAL & paper's finite real remainder $=$ BSV QMRK$-$MRK at $s_\Lambda\to\infty$, Eq.~(\ref{d:schemeid})\\
rapidity scheme, virtual / $s_0$ & not done & split of the trajectory loop at $p^+=k^+$ vs $y_p=y_k$; definition of evolved $\phi_{P,T}$\\
$\ln\mu^2$, scale of $\alpha_s$ & AGREES & $-\frac{\alpha_s\beta_0}{4\pi}\ln\frac{\v k^2}{\mu^2}$ (FL, BSV)\\
collinear region, polarisation scheme & AGREES (amplitude) & BSV $A_{(1)}+A_{(2)}$; CDR structure ($D-2=d_\perp$)\\
$\frac{67}{9}-\frac{\pi^2}{3}$ & OPEN, condition (\ref{d:cond}) & real $=$ FL's real, so Group I must carry no net $c$; Row I.II alone gives $+\frac{\alpha_sN_c}{2\pi}c$; Rows I.III, I.V, contact not evaluated\\
one-loop vertex (Group I) & reduced, NOT EVALUATED & amplitude-level reduction, Eq.~(\ref{d:virtred}); I.II: $-\omega_B(\v K)\ln\frac{k^+}{\Lambda}$, $\v K=\v k,\v q_a$\\
$f^{abc}\{\rho^b,\rho^c\}$ cancellation in Group I & not found & zero for a point projectile; contact terms otherwise\\
printed singlet contact terms vs.\ commutators $C$ & NOT CHECKED & coordinate space after $p^+$ integration and regrouping into $\mathcal S$\\
finite terms as printed (coordinate space) & NOT TESTED & this note tests the amplitudes\\
CXY hybrid formulas, dilute limit & NOT DONE & different kinematics (collinear projectile, forward)\\\hline
\end{tabular}
\end{center}
No disagreement was established. The open points are the evaluation of Group I and the virtual trajectory conversion; without them the constant and the finite kernel are not validated, and condition (\ref{d:cond}) shows that the constant is a genuine test: with Row I.II as the only LO-type constant it would fail.

\subsection{Drop-in text for the paper}
Proposed replacement of the paragraph ``Dilute target'' in Section~\dref{sec:limits} (macros of the draft; labels new). The bibliography keys exist in \texttt{refs.bib}.
\begin{verbatim}
\paragraph{Dilute target and dilute projectile.} For $U=1+u$,
$u^{ab}=f^{acb}\alpha^c$, keeping $O(u)$ in the amplitude and in the
conjugate amplitude and the two-point function of the projectile, the LO
spectrum reduces to the $k_T$-factorised form, Eq.~\eqref{eq:kTfact},
$\frac{d^3N_{\rm LO}}{d^3k}=\mathcal N\int_{\v q}\phi_P(\v q)
\phi_T(\v k-\v q)/\v k^2$, with $\mathcal N=\frac{\alpha_sN_c(N_c^2-1)}
{\pi^2}\frac{S_\perp}{k^+}$, $\phi_P(\v q)=P(\v q)/\v q^2$,
$\phi_T(\v l)=\v l^2T(\v l)$. For the real emission (Groups II and III)
the anticommutator parts of the two-gluon amplitude cancel at $O(u)$,
$\mathcal A_B^{(2)}+\mathcal A^{(1)}_{AA}+\mathcal A^{(2)}_{AA}
=\frac14[L_k,\mf A_p]+\frac14[L_p,\mf A_k]+O(u^2)$, with $L$ the LO
amplitude before the Fourier transform, so that the two-gluon amplitude
is linear in $\rho$, and
\begin{align}\label{eq:dilute_real}
\frac{d^3N^{\rm (II+III)}}{d^3k}\Big|_{\rm dilute}=\mathcal N
\int_{\v q_a}\int_\Lambda^\vee\frac{dp^+}{p^+}\int d^2p\,
\frac{\alpha_sN_c}{\pi^2}\,\phi_P(\v q_a)\phi_T(\v q_b)\,
\frac{A_{gg}(\v q_a,\v q_b;\v k,\v p;x)}{\v q_a^2\v q_b^2}\,,
\end{align}
$\v q_b=\v k+\v p-\v q_a$, $x=k^+/(k^++p^+)$, where $A_{gg}$ is the
Born cross section for the production of two gluons in
quasi-multi-Regge kinematics~\cite{Fadin:1998py}, in the form given in
Ref.~\cite{Bartels:2006hg} (their $A_{\rm gluons}$). We verified this
identity numerically at random kinematics for all $0<x<1$, for $SU(2)$
and $SU(3)$ and for $d_\perp$ transverse polarisations against the
$(D-2)$ terms of~\cite{Bartels:2006hg} (relative deviation
$2\times10^{-15}$, no free normalisation); it requires the commutator
(contact) parts generated by the independent-emission terms. In the
limits $p^+\gg k^+$ and $p^+\ll k^+$, $A_{gg}\to\v q_a^2\v q_b^2/
(\v k^2\v p^2)$, the real parts of one step of LO BFKL evolution of
$\phi_P$ and $\phi_T$ with $\bar\alpha_s=\alpha_sN_c/\pi$; the virtual
parts are those of Sections~\ref{sec:check_BFKL}
and~\ref{sec:check_JIMWLK}. The coefficient of $\ln\mu^2$ and the scale
of the coupling agree with~\cite{Fadin:1998py,Bartels:2006hg}. Since
$dp^+/p^+=dx/(x(1-x))$ is, at fixed $k^+$, the rapidity measure of the
second gluon, the finite part of Eq.~\eqref{eq:dilute_real} (the $x$
integral of $A_{gg}/(\v q_a^2\v q_b^2)-1/(\v k^2\v p^2)$) coincides
identically with the two-gluon contribution minus its multi-Regge limit
of Ref.~\cite{Bartels:2006hg} for $s_\Lambda\to\infty$: for the real
emission the plus-momentum cut-offs require no conversion. The virtual
correction (Group I) in this limit, and its comparison with the
one-loop central-emission
vertex~\cite{Fadin:1998py,DelDuca:1998cx,Bartels:2006hg}, including the
trajectory terms, remain to be done. Because the real emission is that
of~\cite{Fadin:1998py}, and the one-loop vertex contains no
$\frac{67}{9}$, the comparison of the constant
$\frac{67}{9}-\frac{\pi^2}{3}$ is a test of Group I alone.
\end{verbatim}
\noindent The last sentence states a consequence without prejudging the outcome; the authors should evaluate Eq.~(\ref{d:cond}) before the paper states where the second half of the constant resides (Section~\dref{sec:lit_constant}, the \texttt{\textbackslash pending} note there would have to be revised: the ``collinear finite terms'' it refers to are real rows, now fixed to be the QMRK cross section).
In Section~\dref{sec:lit_finite}, item (v), replace the last sentence
\begin{verbatim}
The comparison of the full dilute limit with the jet vertex
of~\cite{Bartels:2006hg,Fadin:1998py} is structural only.
\end{verbatim}
by
\begin{verbatim}
In the dilute limit the real emission agrees exactly with the
quasi-multi-Regge two-gluon production of~\cite{Fadin:1998py,
Bartels:2006hg}, Eq.~\eqref{eq:dilute_real}; the virtual part has not
been compared with the one-loop vertex.
\end{verbatim}

\subsection{What remains open}
\begin{enumerate}
\item The loop integrals of the dilute Group~I, Eq.~(\ref{d:virtred}) (best for a point-like projectile, where Rows~I.I and the $\rho\rho$, $\rho\rho\rho$ parts of $\mf A^{(3)}$ vanish), and the comparison with $K^{(v)}_r$/DDS including the trajectory terms. This decides condition (\ref{d:cond}) for the constant.
\item The virtual trajectory conversion (split of the loop-gluon window at $p^+=k^+$ instead of $y_p=y_k$) and the identification of the cut-off logarithms with BSV's $s_0$, $s_0'$, Sec.~\ref{sec:d-logs}; the real part needs none.
\item The momentum-sum-rule question of the fragmentation function (Major~4), which changes the momentum-weighted constant by $\frac{\alpha_s}{2\pi}N_c(\frac{65}{36}-\frac{2\pi^2}{3})$.
\item Whether the colour-singlet contact terms printed in the draft equal the dilute contact contributions $2{\rm Re}[(\mathcal A_B^{(1)}-\mathcal A_{CA})^\dagger C]+|C|^2$ required by the QMRK check (they are printed in coordinate space, after the $p^+$ integration and the regrouping of the soft remainders into $\mathcal S$; the comparison needs the dilute limit of eight rows and was not done).
\item The dilute limit of the coordinate-space finite terms as printed, against Eq.~(\ref{d:realsplit}) (real) and the result of item 1 (virtual).
\end{enumerate}
